\documentclass[11pt]{article}
\usepackage[a4paper,margin=1in]{geometry}
\usepackage{fontspec}
\usepackage[boldfont=false]{xeCJK}
\setCJKmainfont{FandolSong-Regular.otf}
\usepackage{amsmath}
\usepackage{amssymb}
\usepackage{array}
\usepackage{booktabs}
\usepackage{graphicx}
\usepackage{adjustbox}
\usepackage{enumitem}
\setlist[itemize]{topsep=2pt,partopsep=0pt,parsep=0pt,itemsep=2pt,leftmargin=1.4em}
\setlist[enumerate]{topsep=2pt,partopsep=0pt,parsep=0pt,itemsep=2pt,leftmargin=1.6em}
\usepackage{parskip}
\usepackage{fvextra}
\fvset{breaklines=true,breakanywhere=true,fontsize=\small}
\setlength{\parindent}{0pt}

\begin{document}

\begin{center}
  {\LARGE\bfseries Astronomy and cosmology}\\[0.35em]
  {\large A Level Physics}
\end{center}
\vspace{0.6em}

\section*{Luminosity and radiant flux intensity}

The \textbf{luminosity} 光度 $L$ of a star is the \textbf{total power} 功率 of radiation it gives out — the \textbf{energy} 能量 radiated per second in all directions. Unit: watt (W).

At distance $d$, this power has spread over a sphere of area $4\pi d^{2}$. The \textbf{radiant flux intensity} 辐射通量密度 $F$ (power per unit area) at distance $d$ is

\[
F = \frac{L}{4\pi d^{2}}.
\]

Unit: $\text{W m}^{-2}$. This is the \textbf{inverse-square law} 平方反比定律 for flux: doubling the distance cuts the flux to a quarter. A telescope measures $F$; if $L$ is known, the distance follows:

\[
d = \sqrt{\frac{L}{4\pi F}}.
\]

\par\medskip
\noindent\begin{minipage}{\linewidth}
\centering
\adjustbox{max width=\linewidth, max totalheight=0.42\textheight}{%
  \includegraphics{/Users/joschka/Desktop/Programming/eXercise/output/eXam/handouts/a_level_physics/assets/25-p419-inverse-square-law.pdf}}\\[0.3em]
\emph{The same power spreads over a larger area as distance grows, so flux falls as 1/d²}
\end{minipage}\par\medskip

\section*{Standard candles}

A \textbf{standard candle} 标准烛光 is an object whose \textbf{luminosity is known} from its type. Once you find one in a distant galaxy and measure the flux $F$ from it, you get its distance from $d = \sqrt{L/(4\pi F)}$.

Examples:

\begin{itemize}
\item \textbf{Cepheid variables} 造父变星 (pulsating stars) — the pulsation period is tightly linked to the luminosity, so the period gives $L$.

\item \textbf{Type Ia supernovae} 超新星 — a white dwarf reaching a critical mass and exploding always has about the same peak luminosity.

\end{itemize}

A standard candle gives $L$ without first knowing the distance, so it reaches galaxies far beyond \textbf{parallax} 视差.

\par\medskip
\noindent\begin{minipage}{\linewidth}
\centering
\adjustbox{max width=\linewidth, max totalheight=0.42\textheight}{%
  \includegraphics{/Users/joschka/Desktop/Programming/eXercise/output/eXam/handouts/a_level_physics/assets/25-p420-cepheid-period-luminosity.pdf}}\\[0.3em]
\emph{For Cepheid variables the pulsation period sets the luminosity, making them standard candles}
\end{minipage}\par\medskip

\section*{Stellar surface temperature}

\subsection*{Wien's displacement law}

A hot body gives out a continuous (\textbf{blackbody} 黑体) spectrum with a peak at a \textbf{wavelength} 波长 $\lambda_{\text{max}}$ set by its \textbf{temperature} 温度. \textbf{Wien's displacement law} 维恩位移定律:

\[
\lambda_{\text{max}} T = \text{constant}, \qquad b \approx 2.90 \times 10^{-3}\ \text{m K}.
\]

Hotter stars peak at \textbf{shorter} wavelengths: a cool red star ($\sim 3000\ \text{K}$) peaks in the infrared; the Sun ($\sim 5800\ \text{K}$) peaks near $500\ \text{nm}$; a hot blue-white star ($\sim 20{,}000\ \text{K}$) peaks in the ultraviolet. Measuring $\lambda_{\text{max}}$ gives the surface temperature.

\par\medskip
\noindent\begin{minipage}{\linewidth}
\centering
\adjustbox{max width=\linewidth, max totalheight=0.42\textheight}{%
  \includegraphics{/Users/joschka/Desktop/Programming/eXercise/output/eXam/handouts/a_level_physics/assets/25-p422-blackbody-wien.pdf}}\\[0.3em]
\emph{A hotter black body radiates more, and its peak wavelength shifts towards the blue (Wien's law)}
\end{minipage}\par\medskip

\subsection*{Stefan–Boltzmann law}

A star, treated as a blackbody sphere of radius $r$ and surface temperature $T$, has luminosity

\[
L = 4\pi \sigma r^{2} T^{4},
\]

where $\sigma = 5.67 \times 10^{-8}\ \text{W m}^{-2}\ \text{K}^{-4}$ is the \textbf{Stefan–Boltzmann constant} 斯特藩-玻尔兹曼常量 (the \textbf{Stefan–Boltzmann law} 斯特藩-玻尔兹曼定律). Two strong dependences:

\begin{itemize}
\item $L \propto r^{2}$ — twice the radius, four times the luminosity (same $T$).

\item $L \propto T^{4}$ — twice the temperature, sixteen times the luminosity (same $r$).

\end{itemize}

\subsection*{Estimating a star's radius}

Combine the two laws:

\begin{enumerate}
\item measure $\lambda_{\text{max}}$ → get $T$ from Wien's law.

\item find $L$ (e.g. from flux $F$ and distance $d$: $L = 4\pi d^{2} F$).

\item solve the Stefan–Boltzmann law for $r$: $r = \sqrt{L/(4\pi \sigma T^{4})}$.

\end{enumerate}

This is how astronomers estimate radii of stars they cannot see as a disc.

\section*{Redshift, Hubble's law and the Big Bang}

\subsection*{Cosmological redshift}

The \textbf{spectral lines} 谱线 of light from distant galaxies are seen at \textbf{longer wavelengths} than their known laboratory values — the whole spectrum is stretched towards the red. This is \textbf{redshift} 红移.

\par\medskip
\noindent\begin{minipage}{\linewidth}
\centering
\adjustbox{max width=\linewidth, max totalheight=0.42\textheight}{%
  \includegraphics{/Users/joschka/Desktop/Programming/eXercise/output/eXam/handouts/a_level_physics/assets/25-p424-redshift-spectrum.pdf}}\\[0.3em]
\emph{The hydrogen absorption lines of a distant star are shifted to longer wavelengths — a redshift}
\end{minipage}\par\medskip

Reading it as a Doppler shift, the galaxy is moving \textbf{away}. For $v \ll c$:

\[
\frac{\Delta \lambda}{\lambda} \approx \frac{v}{c},
\]

where $\Delta\lambda = \lambda_{\text{observed}} - \lambda_{\text{emitted}}$ and $v$ is the speed of \textbf{recession} 退行. Example: light emitted at $4.62 \times 10^{-7}\ \text{m}$ but seen at $4.91 \times 10^{-7}\ \text{m}$ gives $\Delta\lambda = 0.29 \times 10^{-7}\ \text{m}$ and

\[
v \approx \frac{\Delta\lambda}{\lambda_{\text{em}}} c \approx 1.9 \times 10^{7}\ \text{m s}^{-1}.
\]

\subsection*{Why redshift means an expanding Universe}

Almost every distant galaxy is redshifted (a few near ones are \textbf{blueshifted} 蓝移 by local motion). So galaxies are, on average, moving apart — not just from us but from each other. The Universe is \textbf{expanding}, with the space between galaxies stretching. More distant galaxies are redshifted more.

\subsection*{Hubble's law}

The link between recession speed $v$ and distance $d$ is \textbf{Hubble's law} 哈勃定律:

\[
v \approx H_{0} \cdot d,
\]

where $H_{0}$ is the \textbf{Hubble constant} 哈勃常数 ($\approx 2.3 \times 10^{-18}\ \text{s}^{-1}$). Always use SI units. Example: a galaxy receding at $1.9 \times 10^{7}\ \text{m s}^{-1}$ is at $d = v/H_{0} \approx 8.3 \times 10^{24}\ \text{m}$.

\subsection*{From Hubble's law to the Big Bang}

Hubble's law means the galaxies were once together. Running the expansion backwards, all distances shrink to zero at $t = -1/H_{0}$ — the Universe was once a tiny, hugely dense, hot point. This is the \textbf{Big Bang} 大爆炸. The age of the Universe (for steady expansion) is about

\[
T_{\text{age}} \approx \frac{1}{H_{0}} \approx 4.3 \times 10^{17}\ \text{s} \approx 14 \text{ billion years}.
\]

The expansion, the redshift of galaxies, the \textbf{cosmic microwave background} 宇宙微波背景, and the hydrogen/helium abundances are the main evidence for the Big Bang.

\subsection*{Distance ladder}

Astronomers combine methods, each calibrated by the one below:

\begin{enumerate}
\item \textbf{parallax} — for nearby stars.

\item \textbf{standard candles} (Cepheids, Type Ia supernovae) — for galaxies.

\item \textbf{Hubble's law} ($d = v/H_{0}$, with $v$ from redshift) — for very distant galaxies.

\end{enumerate}


\end{document}
